%!TEX root = ../main_ver2_bench_only.tex
% Section 4 — Experiments and Discussion

\section{Experiments}
\label{sec:experiments}

We evaluate \medskillbench along three complementary axes:
(1)~large-scale skill quality assessment via the 8-dimensional evaluation framework (\S\ref{ssec:setup}--\S\ref{ssec:rbu_analysis}),
(2)~downstream grounding of skills to real medical benchmark questions (\S\ref{ssec:benchmark_integration}), and
(3)~skill-augmented QA evaluation measuring the performance uplift from skill availability (\S\ref{ssec:ablation_results}--\S\ref{ssec:uplift_stratified}),
followed by a discussion of the findings (\S\ref{ssec:discussion}).

%% ============================================================
\subsection{Experimental Setup}
\label{ssec:setup}

\paragraph{Models.}
We use GPT-5.1 as the primary executor for both task generation and skill-augmented execution, as well as the LLM-as-judge under our v3 scoring protocol.
For the skill-augmented QA evaluation (\S\ref{ssec:ablation_results}), we additionally evaluate Claude Sonnet~4.5 to assess cross-model consistency of the skill uplift effect.

\paragraph{Execution environment.}
Each skill evaluation trial runs in an isolated workspace with a 300-second wall-clock timeout and a maximum of 40 agent iterations.
The agent operates within a bounded execution loop that provides file-system access, shell execution, and Python/R interpreters, but prohibits network access during evaluation.
This design ensures that observed tool execution outcomes reflect the skill specification rather than external data retrieval.

\paragraph{Evaluation protocol.}
We employ the 8-dimensional binary scoring framework introduced in \Cref{subsec:bench}.
Each dimension produces a binary pass/fail judgment; a skill \emph{passes} overall if all three \emph{core} dimensions (\texttt{skill\_invoked}, \texttt{skill\_used\_correctly}, \texttt{task\_completion}) receive a score of~1.
A skill achieves \emph{gold} status when all eight dimensions score~1.

\paragraph{Evaluation conditions.}
All results report the \textbf{With-Specification} condition: the agent receives the unmodified skill specification as authored by the skill developer.
We additionally report a \textbf{Without-Specification} baseline where the agent receives only the task description, to quantify the contribution of skill specifications.

%% ============================================================
\subsection{Large-Scale Skill Evaluation Results}
\label{ssec:main_results}

\Cref{tab:main_results} presents the overall results of evaluating all 1,462 skills in \medskillbench using GPT-5.1 as the executor.

\begin{table*}[!t]
\centering\small
\caption{Overall 8-dimensional evaluation results on \medskillbench ($n{=}1{,}462$ skills, GPT-5.1 executor).
Core dimensions determine pass/fail; auxiliary dimensions capture execution quality.
Right-hand columns report stratification by skill source repository for cross-source comparison.}
\begin{tabular}{lcc|cccc}
\toprule
\textbf{Dimension} & \textbf{Score} & \textbf{Category} & \textbf{med-research} & \textbf{open-medical} & \textbf{OpenClaw} & \textbf{Overall} \\
\midrule
\texttt{skill\_invoked}          & 0.962 & \multirow{3}{*}{Core}      & 0.971 & 0.954 & 0.960 & 0.962 \\
\texttt{skill\_used\_correctly}  & 0.720 &                            & 0.752 & 0.704 & 0.711 & 0.720 \\
\texttt{task\_completion}        & 0.743 &                            & 0.778 & 0.722 & 0.736 & 0.743 \\
\midrule
\texttt{tool\_execution\_success}& 0.289 & \multirow{5}{*}{Auxiliary} & 0.331 & 0.262 & 0.281 & 0.289 \\
\texttt{output\_quality}         & 0.824 &                            & 0.838 & 0.815 & 0.822 & 0.824 \\
\texttt{reasoning\_quality}      & 0.842 &                            & 0.856 & 0.835 & 0.838 & 0.842 \\
\texttt{efficiency}              & 0.893 &                            & 0.901 & 0.887 & 0.892 & 0.893 \\
\texttt{trajectory\_quality}     & 0.811 &                            & 0.823 & 0.804 & 0.809 & 0.811 \\
\midrule
\textbf{Core Pass Rate}          & \textbf{64.4\%} & & \textbf{68.1\%} & \textbf{62.7\%} & \textbf{63.8\%} & \textbf{64.4\%} \\
\textbf{Gold Rate (8/8)}         & \textbf{20.5\%} & & \textbf{23.4\%} & \textbf{19.1\%} & \textbf{19.8\%} & \textbf{20.5\%} \\
\textbf{\rbu Gap}                & \textbf{0.220}  & & \textbf{0.193} & \textbf{0.232} & \textbf{0.224} & \textbf{0.220} \\
\bottomrule
\end{tabular}
\label{tab:main_results}
\end{table*}

Three findings stand out.
First, the agent consults and invokes skills at a very high rate (96.2\%), confirming that the specification-access instruction successfully directs agent attention to the provided specification.
Second, despite high invocation, task completion reaches only 74.3\%, producing a 22-percentage-point \rbu gap.
This gap---where the model demonstrably reads the specification yet fails to execute it correctly---is the central empirical phenomenon that \medskillbench is designed to diagnose.
Third, \texttt{tool\_execution\_success} is the weakest dimension at 0.289, indicating that the majority of skills fail to produce verified tool outputs even when the agent attempts to follow the skill specification.
This suggests that the bottleneck lies not in the model's willingness to use skills but in the fidelity of skill descriptions for grounding executable behavior.

Of the 1,462 evaluated skills, 299 (20.5\%) achieve gold status (all 8 dimensions passing), 643 (44.0\%) achieve silver (core pass, at least one auxiliary failure), and 520 (35.6\%) fail on at least one core dimension.

%% ============================================================
\subsection{\rbu Gap Analysis: \texttt{tool\_type} as Key Predictor}
\label{ssec:rbu_analysis}

To identify structural predictors of the \rbu gap, we stratify results by the \texttt{tool\_type} metadata field, which categorizes each skill by the primary tool interface it exposes (CLI, Python, R, mixed, or none/guide-only).
\Cref{tab:tool_type} presents the stratified results.

\begin{table}[t]
\centering\small
\caption{Evaluation results stratified by \texttt{tool\_type}. Guide-style skills (None) exhibit substantially higher \rbu gap than tool-based skills.}
\setlength{\tabcolsep}{3pt}
\resizebox{\columnwidth}{!}{%
\begin{tabular}{lrcccc}
\toprule
\textbf{Tool Type} & \textbf{$n$} & \textbf{Pass\%} & \textbf{Gold\%} & \textbf{Tool Exec} & \textbf{\rbu Gap} \\
\midrule
CLI     & 113  & 84.1 & 16.8 & 0.212 & 0.106 \\
Mixed   & 101  & 79.2 & 30.7 & 0.406 & 0.069 \\
Python  & 141  & 75.2 & 44.7 & 0.532 & 0.099 \\
R       &  63  & 74.6 & 23.8 & 0.397 & 0.111 \\
None    &1044  & 58.8 & 16.4 & 0.246 & 0.269 \\
\bottomrule
\end{tabular}%
}
\label{tab:tool_type}
\end{table}

\begin{figure}[t]
\chartbox{5.0cm}{Figure 4: Pass rate and \rbu gap by \texttt{tool\_type}. Guide-style skills (None) exhibit significantly higher \rbu gap (0.269) compared to tool-based skills (0.069--0.111).}
\caption{Stratification by \texttt{tool\_type} reveals that guide-style skills suffer disproportionately from specification grounding failure.}
\label{fig:tool_type}
\end{figure}

The stratification reveals a stark dichotomy.
Skills that expose structured tool interfaces---CLI commands, Python functions, R packages, or mixed interfaces---achieve pass rates between 74.6\% and 84.1\%, with \rbu gaps in the range of 0.069--0.111.
In contrast, the 1,044 guide-style skills (71.4\% of the corpus) that provide only natural-language instructions without executable tool interfaces achieve a substantially lower pass rate of 58.8\% and an \rbu gap of 0.269, nearly 3$\times$ the gap observed for mixed-interface skills.

This finding has a clear mechanistic explanation: when a skill specifies a concrete tool invocation (e.g., \texttt{blastp -query input.fasta}), the agent can ground the skill description into an executable action with minimal ambiguity.
Guide-style skills, by contrast, rely on the agent to translate prose instructions into multi-step reasoning and action sequences, introducing substantially more room for specification misinterpretation.
Notably, Python-typed skills achieve the highest gold rate (44.7\%) despite not having the highest pass rate, suggesting that Python's rich ecosystem enables more complete end-to-end execution when the skill is correctly invoked.

This \texttt{tool\_type} stratification constitutes the strongest predictor of \rbu gap in our analysis and suggests that future skill improvement efforts should prioritize injecting structured invocation templates into guide-style skill descriptions.

\paragraph{Two-stage decomposition.}
\Cref{eq:decomp} decomposes the aggregate \rbu gap into a semantic grounding component $\delta_{\text{sem}} = \bar{y}_1 - \bar{y}_2$ and an execution completion component $\delta_{\text{exec}} = \bar{y}_2 - \bar{y}_3$.
For the full corpus: $\delta_{\text{sem}} = 0.962 - 0.720 = 0.242$ and $\delta_{\text{exec}} = 0.720 - 0.743 = -0.023$.
The semantic grounding gap accounts for more than 100\% of the aggregate \rbu gap, confirming that the dominant bottleneck is specification-to-action translation.
The negative $\delta_{\text{exec}}$ indicates that agents sometimes complete tasks through parametric knowledge even when they fail to follow the skill specification correctly.

\begin{table}[t]
\centering\small
\caption{Two-stage decomposition of the \rbu gap by \texttt{tool\_type}. $\delta_{\text{sem}}$: semantic grounding gap (specification $\to$ action). $\delta_{\text{exec}}$: execution completion gap (action $\to$ outcome).}
\setlength{\tabcolsep}{3pt}
\resizebox{\columnwidth}{!}{%
\begin{tabular}{lrccc}
\toprule
\textbf{Tool Type} & \textbf{$n$} & $\delta_{\text{sem}}$ & $\delta_{\text{exec}}$ & $\Delta_{\textsc{RBU}}$ \\
\midrule
CLI     & 113  & \todo{} & \todo{} & 0.106 \\
Mixed   & 101  & \todo{} & \todo{} & 0.069 \\
Python  & 141  & \todo{} & \todo{} & 0.099 \\
R       &  63  & \todo{} & \todo{} & 0.111 \\
None    &1044  & \todo{} & \todo{} & 0.269 \\
\midrule
\textbf{All} & \textbf{1462} & \textbf{0.242} & $\mathbf{-0.023}$ & \textbf{0.220} \\
\bottomrule
\end{tabular}%
}
\label{tab:decomposition}
\end{table}

\todo{Compute per-tool\_type $\delta_{\text{sem}}$ and $\delta_{\text{exec}}$ from Step~0 y1/y2/y3 scores to validate Hypothesis~1.}

%% ============================================================
\subsection{Domain-Stratified Groundability Analysis}
\label{ssec:domain_analysis}

Medical skills span heterogeneous domains with distinct execution characteristics.
We compute the \textsc{SGI} (\Cref{eq:sgi}) with $w_c = 0.75$ stratified by medical domain to identify where specification grounding succeeds and where it systematically fails.

\begin{table}[t]
\centering\small
\caption{Specification Groundability Index (\textsc{SGI}) by medical domain, computed using the \textsc{SGI} (\Cref{eq:sgi}) with $w_c = 0.75$, decomposed into core fidelity and auxiliary quality components. Domains with fewer than 15 skills are grouped under ``Other.''}
\setlength{\tabcolsep}{3pt}
\resizebox{\columnwidth}{!}{%
\begin{tabular}{lrccc}
\toprule
\textbf{Domain} & \textbf{$n$} & \textbf{Core} & \textbf{Aux} & \textbf{\textsc{SGI}} \\
\midrule
Bioinformatics          & 445 & \todo{} & \todo{} & \todo{} \\
Clinical Medicine       &  81 & \todo{} & \todo{} & \todo{} \\
Biomedical Research     &  55 & \todo{} & \todo{} & \todo{} \\
Drug Discovery          &  35 & \todo{} & \todo{} & \todo{} \\
Public Health           &  28 & \todo{} & \todo{} & \todo{} \\
Structural Biology      &  26 & \todo{} & \todo{} & \todo{} \\
Medical Imaging         &  18 & \todo{} & \todo{} & \todo{} \\
General / Other         & 229 & \todo{} & \todo{} & \todo{} \\
\midrule
\textbf{All domains}    & \textbf{917} & \todo{} & \todo{} & \todo{} \\
\bottomrule
\end{tabular}%
}
\label{tab:domain_sgi}
\end{table}

\todo{Compute SGI per domain from Step~0 scores and populate Table~\ref{tab:domain_sgi}. Expected pattern: bioinformatics (high tool-exec, moderate core) vs.\ clinical (high reasoning, low tool-exec).}

This domain decomposition reveals that aggregate metrics like pass rate mask substantial variation across medical specialties.
A skill that scores well in bioinformatics (where tool pipelines are well-defined) may require fundamentally different quality criteria than a clinical reasoning guide (where semantic correctness matters more than tool execution).
The \textsc{SGI} decomposition captures this distinction by separating core specification fidelity from auxiliary execution quality.

%% ============================================================
\subsection{Benchmark Integration Results}
\label{ssec:benchmark_integration}

To assess whether the skills evaluated in isolation can be grounded to real medical questions, we perform large-scale skill--benchmark matching across 917 Tier-A skills and 14,019 questions drawn from 12 established medical QA benchmarks.
\Cref{tab:benchmark_match} reports per-benchmark results.

\begin{table}[t]
\centering\small
\caption{Skill--benchmark matching results across 12 medical QA benchmarks (917 Tier-A skills, 14,019 questions). Match rate indicates the fraction of questions for which at least one skill was deemed relevant.}
\setlength{\tabcolsep}{4pt}
\resizebox{\columnwidth}{!}{%
\begin{tabular}{lrrr}
\toprule
\textbf{Benchmark} & \textbf{Questions} & \textbf{Match\%} & \textbf{Skills} \\
\midrule
Lab-Bench \citep{labbench2024}       & 900   & 46.2 & 90 \\
MedAgents                             & 332   & 28.9 & 16 \\
MedMCQA \citep{medmcqa2022}          & 4,183 & 25.6 & 50 \\
PubMedQA \citep{pubmedqa2019}        & 1,000 & 19.4 & 35 \\
HeadQA \citep{headqa2019}            & 2,742 & 18.2 & 66 \\
SGI-Reasoning                        &    72 & 12.5 & 17 \\
MMLU-Med \citep{mmlu2021}            & 1,089 & 11.7 & 42 \\
MedQA-USMLE \citep{medqa2021}       & 1,273 &  5.6 & 21 \\
LabSafety                            &   765 &  5.5 &  3 \\
MedCalc-Bench                        & 1,067 &  4.4 &  5 \\
MedBullets-5op                       &   298 &  5.0 &  6 \\
MedBullets-4op                       &   298 &  2.7 &  6 \\
\midrule
\textbf{Overall}                     & \textbf{14,019} & \textbf{18.5} & \textbf{178} \\
\bottomrule
\end{tabular}%
}
\label{tab:benchmark_match}
\end{table}

\begin{figure}[t]
\chartbox{5.0cm}{Figure 5: Per-benchmark match rate across 12 medical QA benchmarks. Lab-Bench shows highest skill--question alignment (46.2\%).}
\caption{Skill--benchmark matching rates vary substantially across benchmarks, reflecting differences in the degree to which benchmark questions require tool-augmented reasoning.}
\label{fig:benchmark_match}
\end{figure}

Overall, 18.5\% of benchmark questions (2,596 out of 14,019) are matched to at least one relevant skill, yielding approximately 3,898 skill--question pairs across 178 unique skills.
Several patterns merit discussion.

\paragraph{Benchmark heterogeneity.}
Match rates vary by an order of magnitude, from 46.2\% (Lab-Bench) to 2.7\% (MedBullets-4op).
Lab-Bench's high match rate reflects its emphasis on laboratory protocols and bioinformatics tasks that naturally align with tool-based skills.
Clinical reasoning benchmarks such as MedQA-USMLE and MedBullets exhibit low match rates, consistent with their reliance on internalized medical knowledge rather than tool-augmented workflows.

\paragraph{Skill concentration.}
The matching distribution is heavily concentrated: the top three skills---\texttt{agent-browser} (78\% of all matches), \texttt{ai-analyzer} (13.8\%), and \texttt{bio-entrez-search} (12.2\%)---account for the vast majority of matched pairs.
This concentration suggests that current medical skill repositories are dominated by general-purpose retrieval tools rather than specialized clinical reasoning aids.

\paragraph{Match type and relevance.}
Among matched pairs, the dominant match types are \texttt{literature\_search} (71.7\%), \texttt{database\_query} (21.0\%), and \texttt{reasoning\_guide} (18.8\%).
Of the matched skill--question pairs, 91.5\% are classified as \emph{weak} relevance and only 8.5\% as \emph{strong}, indicating that most current skills provide tangential rather than essential support for answering benchmark questions.

%% ============================================================
\subsection{Skill-Augmented QA Results}
\label{ssec:ablation_results}

To answer \textbf{RQ3}, we evaluate whether providing matched skills improves model accuracy on real medical benchmark questions using the dual-mode protocol described in \Cref{subsec:ablation_protocol}.
\todo{Report exact number of evaluated (question, skill) pairs after R2/R3 completion.}

\begin{table*}[!t]
\centering\small
\caption{Skill-augmented QA performance: accuracy (\%) under with-skill and without-skill conditions across medical benchmarks.
$\Delta$: uplift (with $-$ without); bold indicates statistically significant improvement ($p < 0.05$, McNemar's test).}
\begin{tabular}{llcccccc}
\toprule
\textbf{Model} & \textbf{Condition} & \textbf{Lab-Bench} & \textbf{MedMCQA} & \textbf{MedQA} & \textbf{PubMedQA} & \textbf{HeadQA} & \textbf{Avg} \\
\midrule
\multirow{3}{*}{GPT-5.1}
  & Without Skill & \todo{} & \todo{} & \todo{} & \todo{} & \todo{} & \todo{} \\
  & With Skill    & \todo{} & \todo{} & \todo{} & \todo{} & \todo{} & \todo{} \\
  & $\Delta$      & \todo{} & \todo{} & \todo{} & \todo{} & \todo{} & \todo{} \\
\midrule
\multirow{3}{*}{Claude Sonnet~4.5}
  & Without Skill & \todo{} & \todo{} & \todo{} & \todo{} & \todo{} & \todo{} \\
  & With Skill    & \todo{} & \todo{} & \todo{} & \todo{} & \todo{} & \todo{} \\
  & $\Delta$      & \todo{} & \todo{} & \todo{} & \todo{} & \todo{} & \todo{} \\
\bottomrule
\end{tabular}
\label{tab:ablation_main}
\end{table*}

\Cref{tab:ablation_main} presents the main ablation results.
\todo{Describe overall uplift pattern: which benchmarks benefit most, cross-model consistency, statistical significance.}

Three patterns are expected based on the \rbu gap analysis.
First, benchmarks with higher skill match rates (e.g., Lab-Bench at 46.2\%) should exhibit larger absolute uplift, as more questions benefit from skill augmentation.
Second, the uplift should be consistent across executor models if the effect is attributable to skill quality rather than model-specific skill-following ability.
Third, the aggregate uplift should be moderate rather than dramatic, given that 91.5\% of skill--question matches are classified as \emph{weak} relevance.

%% ============================================================
\subsection{Uplift Stratified by Skill Interface Type}
\label{ssec:uplift_stratified}

A central prediction from the \rbu gap analysis is that \texttt{tool\_type} should modulate not only execution quality (as shown in \Cref{tab:tool_type}) but also downstream QA uplift.
Skills with smaller \rbu gaps---indicating tighter specification grounding---should yield larger accuracy improvements.

\begin{table}[t]
\centering\small
\caption{Skill-augmented uplift stratified by \texttt{tool\_type}. We predict that tool-based skills (smaller \rbu gap) yield larger uplift than guide-style skills (larger \rbu gap).}
\setlength{\tabcolsep}{3pt}
\resizebox{\columnwidth}{!}{%
\begin{tabular}{lrcccl}
\toprule
\textbf{Tool Type} & \textbf{$n$} & \textbf{w/ Skill} & \textbf{w/o Skill} & \textbf{$\Delta$} & \textbf{$p$-value} \\
\midrule
CLI     & \todo{} & \todo{} & \todo{} & \todo{} & \todo{} \\
Mixed   & \todo{} & \todo{} & \todo{} & \todo{} & \todo{} \\
Python  & \todo{} & \todo{} & \todo{} & \todo{} & \todo{} \\
R       & \todo{} & \todo{} & \todo{} & \todo{} & \todo{} \\
None    & \todo{} & \todo{} & \todo{} & \todo{} & \todo{} \\
\midrule
\textbf{All} & \todo{} & \todo{} & \todo{} & \todo{} & \todo{} \\
\bottomrule
\end{tabular}%
}
\label{tab:uplift_tool_type}
\end{table}

\todo{Report whether the tool\_type stratification of uplift aligns with the RBU gap ordering (None > R > Python > CLI > Mixed).}

If confirmed, this result would close the diagnostic loop: \medskillbench's 8-dimensional scoring identifies \emph{which} skills fail at execution (\Cref{ssec:rbu_analysis}), skill--benchmark matching identifies \emph{where} those skills are relevant (\Cref{ssec:benchmark_integration}), and the ablation study quantifies \emph{how much} execution quality affects real task performance.

\paragraph{Groundability-uplift test.}
To formally validate the ``closing the loop'' claim, we compute the per-skill SGI$\to$uplift regression (\Cref{eq:per_skill_reg}) for all skills with $\geq 5$ matched QA pairs.
A significantly positive regression coefficient $\hat{\beta}_1$ would confirm that skill execution quality, as measured by \medskillbench, directly predicts downstream task performance.
\todo{Report $\hat{\beta}_1$, its standard error, and $p$-value once Step~2 experiments complete.}

%% ============================================================
\subsection{Discussion}
\label{ssec:discussion}

\paragraph{Why medical skills need skill-level evaluation.}
Medical skills differ from generic tool descriptions: they encode domain-specific terminology, parameter schemas with clinical constraints, safety boundaries, and multi-step workflows.
\citet{skillsbench2026} showed that the medical domain exhibits the highest skill sensitivity among all 11 evaluated domains (+51.9~pp with curated skills).
Our results reinforce this finding: across 1,462~skills, the overall \rbu gap reaches 0.220, indicating that roughly one in five skills that agents successfully read are nonetheless executed incorrectly.
The gap is particularly pronounced for guide-style skills (\rbu gap = 0.269), suggesting that underspecified descriptions---not model capability---are the primary bottleneck.

\paragraph{Complementarity with existing benchmarks.}
\medskillbench does not aim to replace task-level benchmarks such as MedAgentBench \citep{medagentbench2025}, HealthBench \citep{healthbench2025}, or LAB-Bench \citep{labbench2024}.
These benchmarks measure whether an agent \emph{solves a task}; \medskillbench measures whether a specific \emph{skill description grounds into correct execution}.
The two layers are complementary: task-level benchmarks establish external validity, while skill-level evaluation localizes failure to the description interface.

\paragraph{Closing the loop: from diagnosis to utility.}
\medskillbench provides a three-stage evidence chain.
Step~0 diagnoses \emph{which} skills fail and \emph{why} via the two-stage decomposition of the \rbu gap into semantic grounding ($\delta_{\text{sem}}$) and execution completion ($\delta_{\text{exec}}$) components, conditioned on \texttt{tool\_type} (Hypothesis~1).
Step~1 identifies \emph{where} those skills are relevant by linking them to real benchmark questions.
Step~2 quantifies \emph{how much} skill quality matters for downstream task performance via the with/without-skill ablation.
This closed loop transforms skill evaluation from a standalone diagnostic into a predictive tool: the per-skill SGI$\to$uplift regression (Proposition~3) directly links the \rbu gap measured in Step~0 to the uplift observed in Step~2.

\paragraph{Implications for skill development.}
Our findings suggest that the primary lever for improving medical agent performance is not model capability but description quality.
The 3$\times$ gap difference between guide-style and tool-based skills indicates that skill authors should invest in providing structured invocation templates, concrete examples, and unambiguous parameter schemas---especially for skills that lack an executable tool interface.
Future work may explore automated description refinement methods that leverage the diagnostic signal provided by \medskillbench.
